% Abhakseite „Das kann ich“ – Zone nur im Gesamt (\mitzone)
%% TODO Ich-kann-Sätze: Platzhalter wie in den Aufgabendateien
\begin{abhakseite}
\ifdefined\mitzone
\abhakgruppe{Kennst du schon}
\abhak{1}{Flächeninhalt von Quadrat, Rechteck, Dreieck (Seitenflächen) und Kreis (Grundkreis)}
\abhak{2}{Volumen von Prisma und Zylinder als Grundfläche mal Höhe (der Bezugskörper für das Drittel) und Mantel des Zylinders}
\abhak{3}{Satz des Pythagoras}
\abhak{4}{Quadrieren und hoch drei, Wurzel und Kubikwurzel mit dem Taschenrechner, Rechnen mit π, Runden auf eine sinnvolle Stelle}
\abhak{5}{Ein Drittel und vier Drittel einer Zahl}
\abhak{6}{Formel aus der Formelsammlung entnehmen und nach einer Größe umstellen (Volumenformel mit Bruchfaktor nach der Höhe)}
\abhak{7}{Radius aus Durchmesser}
\abhak{8}{Volumeneinheiten cm³ → dm³ = l mit 1000, m³; Masse aus Volumen und Dichte}
\abhak{9}{Termumformung mit Variablen}
\abhak{10}{Radius aus Durchmesser – Fehler finden}
\abhak{11}{Radius aus Durchmesser – selbst rechnen}
\fi
\abhakgruppe{Einheit 1 · Pyramide}
\abhak{12}{Säule, Spitze oder Kugel?}
\abhak{13}{Stützdreieck einzeichnen}
\abhak{14}{Volumen oder Oberfläche?}
\abhak{15}{Welche Formel?}
\abhak{16}{Pyramide}
\abhak{17}{Pyramide – Fehler finden, Begründen, Darstellungswechsel, Anwendung}
\abhakgruppe{Einheit 2 · Kegel}
\abhak{18}{Radius oder Durchmesser?}
\abhak{19}{Kegel}
\abhak{20}{Kegel – Fehler finden, Begründen, Darstellungswechsel, Anwendung}
\abhakgruppe{Einheit 3 · Kugel}
\abhak{21}{Kugel}
\abhak{22}{Kugel – Fehler finden, Begründen, Darstellungswechsel, Anwendung}
\end{abhakseite}
